Resilient positioning, navigation, and timing (PNT) underpins aviation,
defence, critical infrastructure, and emerging cislunar operations, yet the
field lacks a common basis for comparing how candidate resilience measures --
classical or quantum -- perform when GNSS is denied. Such comparisons are
predominantly private and non-reproducible: the authoritative recent survey of
quantum PNT names no open, reproducible tool for quantum-versus-classical
evaluation \citep{quantumpntreview2026}, and broader GNSS work has documented
how rarely reported results can be regenerated end to end
\citep{fernandezprades2018reproducibility}. We address this with a methodology
in which every result is reproducible from a committed scenario, random seed,
and engine version; every figure of merit carries an explicit
validated-or-not-modeled label; and the quantum/classical axis is a single
parameter swap inside one neutral code path, so the two are never compared
across separate implementations. The methodology is instantiated in Kshana
\citep{baweja2026kshana}, an open engine written in Rust with native, Python,
and WebAssembly targets and 19 scenario kinds. The new result is a Monte-Carlo
crossover/break-even map over the published-parameter envelope, with 95\%
bootstrap confidence intervals, identifying where a cold-atom inertial sensor or
optical-lattice clock changes the GNSS-outage outcome and where it does not.
Trust evidence is reported separately and tiered by strength: reference frames
agree bit-for-bit with SOFA/ERFA, SGP4 matches the AIAA 2006-6753 set to
\SI{4.12}{\milli\metre} worst-case, Allan deviations match NIST references, and
force-model validation by ephemeris fitting against agency precise ephemerides
yields a full-arc Galileo residual of \SI{0.61}{\metre} at degree/order 70.
Quantum-sensor figures are performance-model projections from cited device
literature, not measurements; the spoofing detector is characterised against
published TEXBAT parameters, not raw IQ.
